%% 第7巻　まとめて計算する　本文（第14回編集委員会）
%% 章は \bkinputchap{ch7-NN} で読む（まだない章は、組むときに飛ばして「BKPAGE missing」を出す）
%% 第7巻は巻の色が青なので、部扉の絵の線は白（bkOn）でかく
%% Ver2（全巻にかかる修正事項 No.58）：章の前に Lua のごみ集めをして、メモリの使用量をログに出す（第9巻と同じ）
\providecommand\bkinputchap[1]{\directlua{collectgarbage("collect") texio.write_nl("log","BKMEM #1 "..math.floor(collectgarbage("count")/1024).." MB")}\IfFileExists{#1.tex}{\input{#1}}{\clearpage\typeout{BKPAGE missing #1}}}
\bkinputchap{ch7-00}

%% ======================== 第1部 行列の計算 ========================
\bkPartOpener{1}{第7巻\hspace{0.8em}第1部}{行列の計算}%
  {172/{たくさんの式を、表のまま一度に計算できる？},
   160/{4つの数で、絵をのばしたり回したりできる？},
   148/{100個の連立方程式は、どうやって解く？}}%
  {% 方眼の上の正方形と、行列で変わった平行四辺形（白い線）
   \begin{scope}[bkOn, line width=1.2pt, line cap=round, line join=round]
     \draw[line width=0.6pt] (120,110) -- (172,110);
     \draw[line width=0.6pt] (124,106) -- (124,142);
     \draw (124,110) rectangle (136,122);
     \draw[line width=1.6pt] (124,110) -- (148,116) -- (160,140) -- (136,134) -- cycle;
     \draw[-{Stealth[length=2.6mm]}, line width=1.4pt] (137,125) -- (146,128);
   \end{scope}}%
  {表の形に並べた数を、そのまま1つの「数」のように計算する道具が行列だよ。行列をかけると、点や絵が動き、たくさんの連立方程式が一度に解けるよ。}
\bkPartContents{1}{%% 自動生成（gen.py）
\bktocchap{7-1}{数の表で計算する：行列}{買い物の合計を、表のかけ算で一度に求める}{5}{3}{3}{24}{C}
\bktocchap{7-2}{行列×ベクトル：形を変える機械}{4つの数で、絵をのばしたり回したりできる}{5}{4}{4}{64}{C}
\bktocchap{7-3}{行列×行列：変換をつなぐ}{「回してのばす」と「のばして回す」はちがう}{5}{4}{4}{102}{C}
\bktocchap{7-4}{行列式：面積が何倍になるか}{0になると、平らにつぶれて元にもどせない}{5}{4}{4}{138}{C}
\bktocchap{7-5}{逆行列で連立方程式を解く}{式が100個あっても、コンピュータの手順は同じ}{5}{4}{4}{186}{C}
}%
{行列のたし算・かけ算、行列式、逆行列を、図と数の表の2通りで計算できるようになります。表計算と Python でも確かめます。}
\bkpartbegin{1}{行列の計算}
\bkinputchap{ch7-01}
\bkinputchap{ch7-02}
\bkinputchap{ch7-03}
\bkinputchap{ch7-04}
\bkinputchap{ch7-05}

%% ======================== 第2部 回転と座標の変換 ========================
\bkPartOpener{2}{第7巻\hspace{0.8em}第2部}{回転と座標の変換}%
  {172/{ロボットの手の位置は、どうやって計算する？},
   160/{30°回して45°回すと、何°回したことになる？},
   148/{ドローンの向きは、いくつの数で決まる？}}%
  {% 2つの関節をもつロボットの腕（白い線）
   \begin{scope}[bkOn, line width=1.2pt, line cap=round, line join=round]
     \draw[line width=0.6pt] (118,110) -- (172,110);
     \draw (124,110) -- (130,110) -- (130,113) -- (124,113) -- cycle;
     \draw[line width=2.2pt] (127,113) -- ({127+24*cos(40)},{113+24*sin(40)});
     \draw[line width=2.2pt] ({127+24*cos(40)},{113+24*sin(40)}) -- ({127+24*cos(40)+18*cos(-5)},{113+24*sin(40)+18*sin(-5)});
     \fill (127,113) circle (1.6);
     \fill ({127+24*cos(40)},{113+24*sin(40)}) circle (1.6);
     \draw[line width=0.9pt] (135,113) arc (0:40:8);
   \end{scope}}%
  {回す・向きを変える・座標を取りかえる、をすべて行列のかけ算で書くよ。ロボットの腕やドローンの姿勢、ゲームの画面は、この計算で動いているよ。}
\bkPartContents{2}{%% 自動生成（gen.py）
\bktocchap{7-6}{回転行列：くるっと回す}{30°回して45°回すと、75°回したのと同じ}{4}{4}{4}{238}{C}
\bktocchap{7-7}{3次元の回転と座標の変換}{ドローンの向きを、3つの回転で表す}{6}{5}{5}{274}{C}
}%
{回転行列と同次変換行列で、ロボットの腕の先の位置や、3次元の向きを計算できるようになります。物理シリーズの力学で、そのまま使います。}
\bkpartbegin{2}{回転と座標の変換}
\bkinputchap{ch7-06}
\bkinputchap{ch7-07}

%% ======================== 第3部 固有値と行列の使いみち ========================
\bkPartOpener{3}{第7巻\hspace{0.8em}第3部}{固有値と行列の使いみち}%
  {172/{変換しても、向きが変わらない矢印はある？},
   160/{2つの台車のゆれは、なぜ2つのリズムに分けられる？},
   148/{同じ行列を何回もかけると、どうなる？}}%
  {% のびるだけの向き（固有ベクトル）と、向きの変わる矢印（白い線）
   \begin{scope}[bkOn, line width=1.2pt, line cap=round, line join=round]
     \draw[line width=0.6pt] (118,110) -- (172,110);
     \draw[line width=0.6pt] (124,106) -- (124,142);
     \draw[line width=0.8pt, dash pattern=on 2pt off 2pt] (124,110) -- (156,142);
     \draw[-{Stealth[length=2.8mm]}, line width=1.6pt] (124,110) -- (134,120);
     \draw[-{Stealth[length=2.8mm]}, line width=2.2pt] (124,110) -- (150,136);
     \draw[-{Stealth[length=2.4mm]}, line width=1.0pt] (124,110) -- (138,112);
   \end{scope}}%
  {行列をかけても向きが変わらない、特別な矢印があるよ。その矢印の向きで見ると、こみいった変換も「のびるだけ」に分かれる。ゆれのモードも、くり返しの行方も、ここから見えてくるよ。}
\bkPartContents{3}{%% 自動生成（gen.py）
\bktocchap{7-8}{固有値と固有ベクトル：向きの変わらない方向}{変換しても、向きが変わらない矢印がある}{5}{5}{5}{312}{C}
\bktocchap{7-9}{対角化：まっすぐな方向に分けて考える}{こみいった変換も、のびるだけの変換に分けられる}{5}{5}{5}{354}{C}
\bktocchap{7-10}{行列のくり返し：漸化式とオイラー法を行列で}{何回もかけるとどうなるかは、固有値が教えてくれる}{5}{5}{5}{386}{C}
\bktocchap{7-11}{行列の使いみち}{ロボット、電気回路、画像…行列はどこで働いている？}{5}{3}{5}{430}{R}
}%
{固有値と固有ベクトルを求め、対角化で振動モードやくり返しの計算を見通せるようになります。第6巻の特性方程式とつながり、第9巻の状態方程式の土台になります。}
\bkpartbegin{3}{固有値と行列の使いみち}
\bkinputchap{ch7-08}
\bkinputchap{ch7-09}
\bkinputchap{ch7-10}
\bkinputchap{ch7-11}

%% ======================== 巻末 ========================
\bkinputchap{ch7-back}
